\documentclass[11pt, a4paper]{article}

% --- Packages ---
\usepackage[margin=1in]{geometry}    % Standard 1-inch margins
\usepackage{amsmath, amssymb}        % Math typesetting
\usepackage{graphicx}                % Image handling
\usepackage{hyperref}                % Clickable links
\usepackage{xcolor}                  % Color support
\usepackage{booktabs}                % Professional tables
\usepackage{cite}                    % Citation management
\usepackage{minted}
\usepackage{listings}
\usepackage{booktabs}
\usepackage{float}
\usepackage{multirow}
\usepackage[table]{xcolor}
% --- Link Setup ---
\hypersetup{
    colorlinks=true,
    linkcolor=blue,
    filecolor=magenta,
    urlcolor=cyan,
    citecolor=blue,
}

% --- Title Information ---
\title{\textbf{Project 3 Report}\\ \large INFSCI 2591 - Algorithm Design}
\author{Abhay Kumara Sri Krishna Nandiraju \\ University of Pittsburgh}
\date{Spring 2026}

\begin{document}

\maketitle

% \begin{abstract}
% This is a brief summary of your project. State the problem you are solving, the methodology or architecture you used, and the primary results or performance metrics you achieved. Keep it concise, usually around 150-200 words.
% \end{abstract}

% --- Main Content ---
\section{O1. Theoretical Time Complexity Analysis}

\subsection{A. Time Complexity of Nearest Neighbor Algorithm}

\textbf{Analysis:} 
\begin{itemize}
  \item Let $N$ be the number of nodes in the graph
  \item The outer loop in the algorithm runs for N-1 times to make sure that every node is visited.
  \item The inner loop searches all the N nodes to find the minimum edge weight from the current node to any unvisited node.
  \item Since, outer loop runs $N-1$ times and inner loop runs N times, the overall time complexity of the Nearest Neighbor algorithm is $O(N(N-1))=O(N^2)$.
\end{itemize}

\subsection{B. Time Complexity of Nearest Insertion Algorithm}

\textbf{Analysis:} 
\begin{itemize}
  \item It takes $O(N^2)$ to find the shortest initial edge since it requires checking all pairs of nodes.
  \item The outer while loop runs for N-2 times and inside the loop, finding the closest unvisited node
  requires comparing all remaining unvisited nodes against all nodes. In the worst case, this is $O(N^2)$.
  \item Evaluating the insertion cost at every possible best position in the tour takes $O(N)$ time.
  \item Overall time complexity is $O((N-2)(N^2 + N)) = O(N^3-2N^2 +N^2 +2N) = O(N^3)$
\end{itemize}


\section{O2. Results Table}

\begin{table}[htbp]
\centering
\renewcommand{\arraystretch}{1.2}
% This resizebox forces the table to fit the width of your text so it won't get cut off
\resizebox{\textwidth}{!}{%
\begin{tabular}{|c|c|l|c|}
\hline
\textbf{Algorithm} & \textbf{File Name} & \textbf{Tour: Sequence of Nodes} & \textbf{Total Weight} \\ \hline
\multirow{9}{*}{Nearest Neighbor} 
& 4n.csv & 0,1,2,3,0 & 2611.0 \\ \cline{2-4}
& 5n.csv & 0,3,2,4,1,0 & 1290.0 \\ \cline{2-4}
& 6n.csv & 0,5,3,4,1,2,0 & 3198.0 \\ \cline{2-4}
& 7n.csv & 0,6,4,3,1,5,2,0 & 2425.0 \\ \cline{2-4}
& 8n.csv & 0,2,4,1,6,3,5,7,0 & 2898.0 \\ \cline{2-4}
& 9n.csv & 0,8,1,6,5,2,3,7,4,0 & 2842.0 \\ \cline{2-4}
& 10n.csv & 0,7,8,4,2,6,5,3,1,9,0 & 3008.0 \\ \cline{2-4}
& 11n.csv & 0,1,6,4,8,2,9,7,5,10,3,0 & 3350.0 \\ \cline{2-4} 
& 12n.csv & 0,8,3,5,6,4,7,11,10,2,1,9,0 & 2646.0 \\ \cline{2-4}
& 13n.csv & 0,9,2,11,7,12,3,6,1,5,4,8,10,0 & 3653.0 \\ \hline
\multirow{9}{*}{Nearest Insertion} 
& 4n.csv & 0,3,2,1,0& 2611.0 \\ \cline{2-4}
& 5n.csv & 0,3,2,4,1,0 & 1290.0 \\ \cline{2-4}
& 6n.csv &0,5,3,2,4,1,0& 2455.0 \\ \cline{2-4}
& 7n.csv &0,1,5,2,3,4,6,0& 1993.0  \\ \cline{2-4}
& 8n.csv & 0,2,4,5,3,7,1,6,0 & 2692.0 \\ \cline{2-4}
& 9n.csv & 0,8,4,2,5,6,1,7,3,0 & 2588.0 \\ \cline{2-4}
& 10n.csv &0,9,1,3,8,4,5,6,2,7,0 & 3423.0  \\ \cline{2-4}
& 11n.csv & 0,3,5,2,9,7,8,4,10,6,1,0 & 2867.0 \\ \cline{2-4}
& 12n.csv &0,4,7,11,10,6,5,1,2,3,8,9,0 & 2930.0  \\ \cline{2-4}
& 13n.csv &0,11,8,3,6,4,5,1,10,7,12,2,9,0 & 3031.0  \\ \hline
\multirow{9}{*}{Dynamic Programming} 
& 4n.csv & 0,1,2,3,0 & 2611.0 \\ \cline{2-4}
& 5n.csv & 0,1,2,4,3,0 & 1146.0 \\ \cline{2-4}
& 6n.csv &0,1,4,2,3,5,0 & 2455.0 \\ \cline{2-4}
& 7n.csv &0,1,5,2,3,4,6,0 & 1993.0  \\ \cline{2-4}
& 8n.csv & 0,2,4,1,7,5,3,6,0 & 2667.0 \\ \cline{2-4}
& 9n.csv & 0,3,2,4,1,6,5,7,8,0 & 2439.0 \\ \cline{2-4}
& 10n.csv &0,5,6,2,4,9,1,3,8,7,0 & 2953.0  \\ \cline{2-4}
& 11n.csv & 0,1,6,10,5,2,9,7,8,4,3,0 & 2701.0 \\ \cline{2-4}
& 12n.csv &0,6,4,7,5,3,8,1,2,10,11,9,0 & 1907.0  \\ \cline{2-4}
& 13n.csv &0,11,2,9,10,7,8,4,5,1,6,3,12,0 & 2523.0  \\ \hline
\multirow{9}{*}{Branch and Bound} 
& 4n.csv & 0,1,2,3,0 & 2611.0 \\ \cline{2-4}
& 5n.csv & 0,1,2,4,3,0 & 1146.0 \\ \cline{2-4}
& 6n.csv &0,1,4,2,3,5,0 & 2455.0 \\ \cline{2-4}
& 7n.csv & 0,1,5,2,3,4,6,0 & 1993.0  \\ \cline{2-4}
& 8n.csv & 0,2,4,1,7,5,3,6,0 & 2667.0 \\ \cline{2-4}
& 9n.csv & 0,3,2,4,1,6,5,7,8,0 & 2439.0 \\ \cline{2-4}
& 10n.csv &0,5,6,2,4,9,1,3,8,7,0 & 2953.0  \\ \cline{2-4}
& 11n.csv & 0,1,6,10,5,2,9,7,8,4,3,0 & 2701.0 \\ \cline{2-4}
& 12n.csv &0,6,4,7,5,3,8,1,2,10,11,9,0 & 1907.0  \\ \cline{2-4}
& 13n.csv &0,11,2,9,10,7,8,4,5,1,6,3,12,0 & 2523.0  \\ \hline
\end{tabular}%
}
\caption{Test Case Results for all 4 algorithms}
\label{tab:tsp_results}
\end{table}
\end{document}
